\documentclass[11pt,letterpaper]{article}
\usepackage[letterpaper,top=0.85in,bottom=0.85in,left=1in,right=1in]{geometry}
\usepackage[T1]{fontenc}
\usepackage{mathptmx}
\usepackage{microtype}
\usepackage[hidelinks]{hyperref}
\usepackage{titlesec}
\setlength{\parindent}{0pt}
\setlength{\parskip}{4pt}
\titleformat{\section}{\normalfont\normalsize\bfseries}{}{0pt}{}
\titlespacing*{\section}{0pt}{7pt}{2pt}
\newcommand{\rg}[1]{\par\smallskip{\normalfont\bfseries\itshape #1}\par}
\newcommand{\task}[1]{\textbf{#1}\hspace{0.4em}}
\hypersetup{pdftitle={Trustworthy Collective Reasoning and Evaluation for AI Agents},pdfauthor={Xinya Du}}
\begin{document}
\begin{center}
{\fontsize{14}{16}\selectfont\bfseries Trustworthy Collective Reasoning and Evaluation for AI Agents\par}
\vspace{2pt}
{\small Xinya Du | Assistant Professor of Computer Science | The University of Texas at Dallas | xinya.du@utdallas.edu\par}
{\small ONR FY27 YIP | Program Officer: Dr.\ Predrag Neskovic, Code 311, Mathematical Data Science\par}
{\small Focus area: AI and Autonomy (Intelligent and Collaborative Agents; Adversarial AI and Autonomy)\par}
\end{center}

\section*{1\quad Problem}
Five AI agents are asked whether a cargo vessel has left port. Four answer yes and one answers no, so a majority vote favors yes, four to one. Yet each yes answer cites a summary that traces back to the same morning report, while the dissenting agent cites a later port-authority update. The evidence is one report against one newer record, not four against one. Intelligence analysts call this \emph{circular reporting}: several voices that appear independent but share a single origin.

Language-model agents now routinely evaluate, critique, and debate one another's answers. Their agreement can carry far less evidence than it appears to, and interaction makes this pressing: discussion lets agents correct one another, but it also makes their judgments dependent. \textbf{The central problem of this project is to determine how much a group of AI agents' agreement is worth as evidence, and to obtain more evidence when agreement is not enough.}

Agreement misleads in three ways. \emph{Shared mistakes:} agents rely on the same source or share the same blind spot. \emph{Self-favoring and strategic ratings:} an agent rates its own answer, or answers like its own, too highly, and could lower competitors' ratings to raise its own standing. \emph{Influence without evidence:} during discussion, agents may give way to whoever is most confident or persistent, including agents that argue dishonestly or coordinate with one another. Research goal 1 (RG1) measures what agreement is worth despite shared mistakes and strategic ratings; RG2 obtains more evidence by checking claims or by debate; and RG3 trains agents so that their agreement carries more evidence. Sharing a source does not by itself make two judgments redundant, because agents can reason about the same source differently, so we estimate dependence from how often agents' errors coincide rather than assume it from overlapping inputs. We focus on factual claims and reasoning steps that can be checked, at least in part.

\section*{2\quad Why it is important}
Using one model to judge another is now standard practice in AI evaluation and oversight \cite{zheng2023judging,verga2024juries}, and multi-agent systems increasingly decide by vote or discussion. Yet errors are strongly correlated across language models \cite{kim2025correlated}, so a panel of several judges may supply only a few effectively independent votes \cite{kohli2026nine}. If agreement is miscounted, every evaluation and decision built on it inherits the error.

Naval information settings sharpen each part of the problem. Reports from different units or sensors often trace to a common origin. Ground truth arrives late and sparsely, if at all. An adversary may inject false reports or compromise an agent. Checking a claim can be costly, or impossible under limited connectivity. And a human decision-maker must judge whether an AI consensus deserves trust. The same question bears on AI-assisted peer review and on overseeing AI systems more capable than their evaluators.

\section*{3\quad Current approaches and their limitations}
\textbf{LLM judges and juries} make evaluation scalable \cite{zheng2023judging,verga2024juries}, but counting votes treats correlated errors as independent, and judges tend to favor their own answers. \textbf{Multi-agent debate} lets models critique and revise one another \cite{du2024debate}, but models also abandon correct answers under social pressure \cite{sharma2024sycophancy,laban2023flipflop}, and when several agents repeat one view, plain debate gives that view several voices. \textbf{AI safety debate} places two debaters before a judge \cite{irving2018debate}; restricting debaters to verified quotations helps judges reach the truth \cite{khan2024debating}. These protocols assume two sides and a separate judge, and do not model the reliability or dependence of many agents.

\textbf{Reliability estimation without labels} has a long statistical history: estimating raters' error rates without ground truth \cite{dawid1979maximum}, jointly with item difficulty \cite{whitehill2009whose}, detecting copying among data sources \cite{dong2009dependence}, and allowing some dependence among classifiers \cite{jaffe2016dependent}. These methods treat raters as black boxes whose judgments never change, and infer dependence only from patterns of votes. Language-model agents instead expose their sources and reasoning, and they change one another's judgments through interaction.

\textbf{Computational social choice} offers principled guarantees for aggregation and selection. Impartial peer review ensures that a participant's ratings of others cannot change its own chance of selection \cite{kurokawa2015impartial,alon2011sum}. Related work addresses fair and manipulation-robust panel selection \cite{flanigan2021fair,flanigan2024manipulation}, when delegating votes to better-informed participants tracks the truth \cite{berinsky2025liquid}, robust aggregation of AI evaluations \cite{khalaf2026maximal}, and manipulation of reviewer assignment in peer review \cite{jecmen2020manipulation}. These mechanisms provide fairness and incentive guarantees for judgments taken as given. They do not model where judgments come from, how discussion changes them, or whether the claims being judged can be checked. We adopt their impartiality guarantees (Task 1.2) and add these elements.

\textbf{The gap:} existing methods either aggregate fixed judgments or let agents interact; none measures how much independent evidence agreement contains, or acquires more evidence when it falls short.

\section*{4\quad Proposed research: what is new, and why it will succeed}
\textbf{What is new.} (i) Agreement is valued by the measured dependence among agents' errors, using the sources and reasoning that language-model agents expose, rather than by counting votes. (ii) Aggregation weights agents by demonstrated skill yet remains impartial, so no agent can raise its own standing through its ratings of others. (iii) A debate protocol for many agents of correlated reliability lets only verifiable evidence change the answer, and agents are trained to contribute evidence rather than to win.

\textbf{Why it will succeed.} The PI's group built the components this work combines. PRD introduced peer ranking and peer discussion among language-model evaluators, achieving higher accuracy and closer agreement with human judgments than relying on the strongest single model, which favors its own answers and is sensitive to answer order \cite{li2024prd}; our PeerEval leaderboard (peereval-ai.github.io) applies peer evaluation across many current models. FG-PRM detects six types of reasoning errors at the level of individual steps \cite{li2025fgprm}. HiPRAG trains search agents with step-level rewards to retrieve only when needed, reaching 65.4\% and 67.2\% average accuracy across seven question-answering benchmarks with 3B and 7B models while cutting unnecessary searches to 2.3\% \cite{wu2026hiprag}. MEQA and the PI's event-extraction research supply tools for generating multi-hop reasoning data \cite{li2024meqa}. The PI's 2024 NSF CAREER award and 2023 Amazon Research Award on process-guided fine-tuning provide foundations for this work, and as a postdoctoral researcher the PI led system integration for the RESIN-11 system in DARPA KAIROS.

\textbf{Setting.} Each agent first answers independently, splitting its reasoning into checkable claims and citing its sources; this sealed first answer is kept as a reference. Agents then exchange messages or debate, and may revise. The unknowns are whether each claim is correct, each agent's skill on each type of task, and how agents' errors depend on one another. The information available is a few audited claims of known correctness, and checks (retrieving a record, testing a counterexample, running a verifier) that have a cost. The goal is accurate, calibrated decisions within a checking budget that remain robust when some agents are dishonest, and that report ``unresolved'' when the evidence is insufficient.

\rg{RG1. Measure what agreement is worth}
\task{Task 1.1: Dependence-aware reliability.} A simple baseline treats endorsements that rest on the same source as one cluster that counts once. Because shared sources need not produce shared conclusions, we will instead estimate dependence from how often agents' errors coincide. Correlated errors shrink agreement to an effective number of independent judgments: $n$ judges whose errors have average correlation $\rho$ count as about $n/(1+(n-1)\rho)$, so five judges with $\rho = 0.5$ count as fewer than two. We will fit a statistical model, in the tradition of Dawid and Skene, that estimates each agent's skill on each task type, each item's difficulty, and shared error factors whose presence is predicted by observable overlap: shared sources, the same underlying model, and overlapping reasoning steps. Clusters are then formed from estimated error correlation, not from shared inputs alone. When errors are independent, agreement alone can reveal each agent's accuracy; when errors are shared, in general it cannot. We will therefore measure how many audited claims are needed to distinguish an accurate cluster from one that repeats a mistake, and test whether the method recovers the dependence planted in our testbeds. PRD and PeerEval judgments provide baselines and data.

\task{Task 1.2: Impartial, skill-weighted aggregation.} Weighting agents by skill can conflict with impartiality. In PRD's ranking, an agent that lowers a rival's score also lowers that rival's weight as a judge, which can raise the agent's own score. Impartial mechanisms prevent this by design but weight all reviewers equally \cite{kurokawa2015impartial}. We will derive each agent's weight as a judge only from audited claims, which it can improve only by being right. Following the partition principle of impartial selection \cite{alon2011sum}, ratings will flow only between groups, and each group's outcomes will be decided by the other groups' ratings, so an agent's ratings never affect its own outcome. Agents built on the same underlying model, or sharing sources, will fall in one group, which also removes self- and family-favoring. For selecting top candidates, as on a leaderboard, we will adapt randomized impartial selection \cite{kurokawa2015impartial} and measure how much accuracy impartiality costs. Agreement within a coalition that the estimated dependence does not explain will be flagged as possible collusion. Language-model reviewers make it practical to test impartial mechanisms at scale, including in scientific peer review.

\rg{RG2. Obtain more evidence: check when possible, debate when not}
Both tasks follow one rule: \emph{a challenge can change the group's answer only through verifiable evidence}, meaning a record or test result obtained by a check, a quotation verified against its source, or a reasoning step reduced to one simple enough to verify directly. Confidence, persistence, and the number of voices do not count.

\task{Task 2.1: Find the disagreement that matters, and check it.} Independently produced solutions are split into claims and aligned, and the method locates where they diverge. A divergence is \emph{consequential} if reversing the disputed claim would change the final answer; step-level error detection from FG-PRM \cite{li2025fgprm} ranks the suspect steps. In the vessel example, the crux is whether the later port update concerns the same vessel on the same day. The method then selects a check suited to the claim: a record with independent provenance, a counterexample, or an executable verifier. Checks are prioritized by how likely they are to change the answer relative to their cost.

\task{Task 2.2: Evidence-grounded debate among many agents.} Some disputes cannot be settled by a single check: the solutions differ too much to align, or no record, tool, or connection is available to test the disputed claim directly. Even then, quotations can be verified against their sources, and arguments can be broken into steps simple enough to verify. For these disputes we will develop a debate protocol for many agents of uncertain and correlated reliability. (i) Each estimated dependence cluster from RG1 sends one representative, so four agents echoing one report receive one voice rather than four. (ii) Debaters question each other's arguments until the disagreement narrows to a single claim \cite{irving2018debate}; if that claim becomes checkable, Task 2.1 checks it. (iii) Only verified quotations and directly verifiable steps count as evidence \cite{khan2024debating}. (iv) A held-out panel of agents that took no side serves as judge, weighted by RG1's reliability estimates; when every agent has taken a side, a separately designated judge model decides, and no agent judges a debate involving its own answer. The outcome is an answer with reduced confidence, or ``unresolved'' with a record of what remains disputed. To expose dishonest agents, we will track three behavioral signals, namely arguing against one's own sealed first answer without new evidence, citing evidence that does not support the claim, and changing others' answers without evidence, and feed them into RG1's weights. We will test when debate adds information beyond voting and how many dishonest debaters the protocol tolerates.

\rg{RG3. Learn from verified outcomes}
\task{Task 3.1: Close the loop.} Each dispute settled by a check becomes a new audited claim that updates RG1's estimates of skill and dependence. Because these claims are selected for being disputed, a small random sample of audited claims is maintained to keep the estimates unbiased. Outcomes decided by debate update the estimates only provisionally and never serve as audits, so persuasive errors cannot reinforce themselves. We will also learn which checks and debate formats resolve which kinds of dispute, and when checking is worth its cost, extending HiPRAG's rewards for searching only when needed \cite{wu2026hiprag}.

\task{Task 3.2: Train agents to contribute, not to win.} Training debaters to win, and using more persuasive debaters, can improve judge accuracy \cite{arnesen2024selfplay,khan2024debating}. These results concern two debaters and a neutral judge; with many correlated and possibly colluding agents, winning rewards persuasiveness that need not coincide with honesty. We will train agents with reinforcement learning on multi-agent and debate episodes whose rewards come from verified outcomes, building on multi-agent training of language models \cite{park2025maporl,subramaniam2025multiagent}. The difficulty is credit assignment: deciding which agent's message moved the group to a correct answer. Building on FG-PRM's step-level rewards and the PI's process-guided fine-tuning, each message will be rewarded by its effect. Introducing verified evidence that moves the group toward the correct answer, or revising one's own answer in response to such evidence, earns reward. Changing others' answers without evidence, giving way without evidence, or abandoning a correct first answer is penalized. The key experiment compares training to win with training to contribute, measuring accuracy, harmful revisions, and emergent dishonest persuasion.

\rg{Task 4 (shared). Testbeds with known dependence, and evaluation}
All scenarios are synthetic or drawn from public data. (A) Math and code problems in which a subset of agents receives the same flawed hint, so dependence is known by construction. (B) Report collections generated from \emph{source trees}, in which each report is an original or a copy, paraphrase, or outdated version of another, using the question-generation pipeline behind our MEQA benchmark \cite{li2024meqa}; these include maritime scenarios with planted circular reporting, alongside MEQA and the multi-hop benchmark MuSiQue \cite{trivedi2022musique}. (C) Scientific claims and peer review: whether a paper's evidence supports its claim, using public conference reviews and author rebuttals, with claims later confirmed or contradicted by reproducibility studies serving as audits. Agent pools will span several model families and include planted dishonest or colluding agents, created by prompting or by training. Baselines are majority vote, a single strong judge, Dawid--Skene-style weighted voting, PRD, impartial selection, plain multi-agent debate, and two-sided debate with verified quotations, all under matched compute and checking budgets. Outcomes are accuracy, calibration, how often a correct minority prevails, how often a correct answer is changed to a wrong one, recovery of the planted dependence, audits needed, the accuracy cost of impartiality, robustness to dishonest agents, and cost. Reference answers used for scoring are kept separate from the audits available to the agents.

\section*{5\quad Limitations}
If every agent shares the same misconception and nothing can be checked, neither aggregation nor debate can recover the truth; our methods should recognize this and report the claim as unresolved. Checks can themselves be wrong, and debate judges can be persuaded, so we will measure sensitivity to unreliable checks and judges rather than treat either as ground truth. Robustness will be established against the dishonest and colluding behaviors we test, not against arbitrary collusion. Reliability estimates may transfer poorly to new topics or model families. Operational naval use would require domain data and expert validation beyond this fundamental research.

\section*{6\quad Expected impact}
If successful, the project will establish when agreement among AI agents counts as evidence and how much checking is needed when it does not, together with impartial skill-weighted aggregation, a many-agent debate protocol in which only verifiable evidence changes the answer, and training that rewards contribution over winning. The naval relevance follows from the setting: circular reporting (Tasks 1.1 and 2.2), late and sparse ground truth (Tasks 1.1 and 3.1), adversarial deception or a compromised agent (Tasks 1.2, 2.2, and 3.2), and costly or unavailable checks (Tasks 2.1 and 2.2). For an analyst combining reports, the resulting capability would trace an assessment to its independent evidence, identify the claim still in dispute, and explain any change of answer or request for review. We will release code, testbeds, and evaluation protocols.

\clearpage
\begingroup
\small
\setlength{\parskip}{0pt}
\begin{thebibliography}{99}
\setlength{\itemsep}{3pt}
\bibitem{zheng2023judging} L. Zheng et al. (2023). \href{https://proceedings.neurips.cc/paper_files/paper/2023/hash/91f18a1287b398d378ef22505bf41832-Abstract.html}{Judging LLM-as-a-Judge with MT-Bench and Chatbot Arena}. NeurIPS 36, Datasets and Benchmarks, pp. 46595--46623.
\bibitem{verga2024juries} P. Verga et al. (2024). \href{https://arxiv.org/abs/2404.18796}{Replacing Judges with Juries: Evaluating LLM Generations with a Panel of Diverse Models}. arXiv:2404.18796.
\bibitem{kim2025correlated} E. M. Kim, A. Garg, K. Peng, and N. Garg (2025). \href{https://proceedings.mlr.press/v267/kim25e.html}{Correlated Errors in Large Language Models}. ICML, PMLR 267:30038--30066.
\bibitem{kohli2026nine} G. Kohli (2026). \href{https://arxiv.org/abs/2605.29800}{Nine Judges, Two Effective Votes: Correlated Errors Undermine LLM Evaluation Panels}. arXiv:2605.29800.
\bibitem{du2024debate} Y. Du, S. Li, A. Torralba, J. B. Tenenbaum, and I. Mordatch (2024). \href{https://proceedings.mlr.press/v235/du24e.html}{Improving Factuality and Reasoning in Language Models through Multiagent Debate}. ICML, PMLR 235:11733--11763.
\bibitem{sharma2024sycophancy} M. Sharma et al. (2024). \href{https://proceedings.iclr.cc/paper_files/paper/2024/hash/0105f7972202c1d4fb817da9f21a9663-Abstract-Conference.html}{Towards Understanding Sycophancy in Language Models}. ICLR.
\bibitem{laban2023flipflop} P. Laban, L. Murakhovs'ka, C. Xiong, and C.-S. Wu (2023). \href{https://arxiv.org/abs/2311.08596}{Are You Sure? Challenging LLMs Leads to Performance Drops in the FlipFlop Experiment}. arXiv:2311.08596.
\bibitem{irving2018debate} G. Irving, P. Christiano, and D. Amodei (2018). \href{https://arxiv.org/abs/1805.00899}{AI Safety via Debate}. arXiv:1805.00899.
\bibitem{khan2024debating} A. Khan et al. (2024). \href{https://proceedings.mlr.press/v235/khan24a.html}{Debating with More Persuasive LLMs Leads to More Truthful Answers}. ICML, PMLR 235:23662--23733.
\bibitem{dawid1979maximum} A. P. Dawid and A. M. Skene (1979). \href{https://doi.org/10.2307/2346806}{Maximum Likelihood Estimation of Observer Error-Rates Using the EM Algorithm}. Journal of the Royal Statistical Society, Series C (Applied Statistics) 28(1):20--28.
\bibitem{whitehill2009whose} J. Whitehill, P. Ruvolo, T. Wu, J. Bergsma, and J. Movellan (2009). \href{https://proceedings.neurips.cc/paper/2009/hash/f899139df5e1059396431415e770c6dd-Abstract.html}{Whose Vote Should Count More: Optimal Integration of Labels from Labelers of Unknown Expertise}. NeurIPS 22, pp. 2035--2043.
\bibitem{dong2009dependence} X. L. Dong, L. Berti-Equille, and D. Srivastava (2009). \href{https://doi.org/10.14778/1687627.1687689}{Integrating Conflicting Data: The Role of Source Dependence}. PVLDB 2(1):550--561.
\bibitem{jaffe2016dependent} A. Jaffe, E. Fetaya, B. Nadler, T. Jiang, and Y. Kluger (2016). \href{https://proceedings.mlr.press/v51/jaffe16.html}{Unsupervised Ensemble Learning with Dependent Classifiers}. AISTATS, PMLR 51:351--360.
\bibitem{kurokawa2015impartial} D. Kurokawa, O. Lev, J. Morgenstern, and A. D. Procaccia (2015). \href{https://procaccia.info/wp-content/uploads/2020/03/impartial.pdf}{Impartial Peer Review}. IJCAI, pp. 582--588.
\bibitem{alon2011sum} N. Alon, F. Fischer, A. D. Procaccia, and M. Tennenholtz (2011). \href{https://doi.org/10.1145/2000378.2000390}{Sum of Us: Strategyproof Selection from the Selectors}. TARK, pp. 101--110.
\bibitem{jecmen2020manipulation} S. Jecmen et al. (2020). \href{https://neurips.cc/virtual/2020/poster/18898}{Mitigating Manipulation in Peer Review via Randomized Reviewer Assignments}. NeurIPS 33.
\bibitem{flanigan2021fair} B. Flanigan, P. G\"olz, A. Gupta, B. Hennig, and A. D. Procaccia (2021). \href{https://doi.org/10.1038/s41586-021-03788-6}{Fair Algorithms for Selecting Citizens' Assemblies}. Nature 596:548--552.
\bibitem{flanigan2024manipulation} B. Flanigan, J. Liang, A. D. Procaccia, and S. Wang (2024). \href{https://procaccia.info/wp-content/uploads/2023/08/sortmanip.pdf}{Manipulation-Robust Selection of Citizens' Assemblies}. AAAI 38(9):9696--9703.
\bibitem{berinsky2025liquid} A. J. Berinsky, D. Halpern, J. Y. Halpern, A. Jadbabaie, E. Mossel, A. D. Procaccia, and M. Revel (2025). \href{https://doi.org/10.1287/mnsc.2023.02470}{Tracking Truth with Liquid Democracy}. Management Science 71(9):7821--7839.
\bibitem{khalaf2026maximal} H. Khalaf, S. L. Wang, D. Halpern, I. Shapira, F. P. Calmon, and A. D. Procaccia (2026). \href{https://arxiv.org/abs/2602.21297}{Robust AI Evaluation through Maximal Lotteries}. ICML.
\bibitem{li2024prd} R. Li, T. Patel, and X. Du (2024). \href{https://openreview.net/forum?id=YVD1QqWRaj}{PRD: Peer Rank and Discussion Improve Large Language Model based Evaluations}. TMLR.
\bibitem{li2025fgprm} R. Li, Z. Luo, and X. Du (2025). \href{https://aclanthology.org/2025.findings-emnlp.228/}{FG-PRM: Fine-grained Hallucination Detection and Mitigation in Language Model Mathematical Reasoning}. Findings of EMNLP, pp. 4247--4278.
\bibitem{wu2026hiprag} P. Wu, M. Zhang, K. Wan, W. Zhao, K. He, X. Du, and Z. Chen (2026). \href{https://proceedings.iclr.cc/paper_files/paper/2026/hash/85255cf8f2d2318cc7cffd0b1bc66625-Abstract-Conference.html}{HiPRAG: Hierarchical Process Rewards for Efficient Agentic Retrieval Augmented Generation}. ICLR.
\bibitem{li2024meqa} R. Li, Z. Wang, S. Q. Tran, L. Xia, and X. Du (2024). \href{https://proceedings.neurips.cc/paper_files/paper/2024/hash/e560a0b22e4432003d0dba63ff8dc457-Abstract.html}{MEQA: A Benchmark for Multi-hop Event-centric Question Answering with Explanations}. NeurIPS 37, Datasets and Benchmarks, pp. 126835--126862.
\bibitem{trivedi2022musique} H. Trivedi, N. Balasubramanian, T. Khot, and A. Sabharwal (2022). MuSiQue: Multihop Questions via Single-hop Question Composition. TACL 10:539--554.
\bibitem{arnesen2024selfplay} S. Arnesen, D. Rein, and J. Michael (2024). \href{https://arxiv.org/abs/2409.16636}{Training Language Models to Win Debates with Self-Play Improves Judge Accuracy}. arXiv:2409.16636.
\bibitem{park2025maporl} C. Park, S. Han, X. Guo, A. E. Ozdaglar, K. Zhang, and J.-K. Kim (2025). \href{https://aclanthology.org/2025.acl-long.1459}{MAPoRL: Multi-Agent Post-Co-Training for Collaborative Large Language Models with Reinforcement Learning}. ACL, pp. 30215--30248.
\bibitem{subramaniam2025multiagent} V. Subramaniam, Y. Du, J. B. Tenenbaum, A. Torralba, S. Li, and I. Mordatch (2025). \href{https://proceedings.iclr.cc/paper_files/paper/2025/hash/1d50221b03ced7152d8750c9d1e66cba-Abstract-Conference.html}{Multiagent Finetuning: Self Improvement with Diverse Reasoning Chains}. ICLR.
\end{thebibliography}
\endgroup
\end{document}
